\subsection{Why conjugate local generators do not suffice}

Nonzero first-cohomology classes with locally varying coboundary parameters
are standard local--global cohomology phenomena; see
\citet[\S2 and Examples~3.4--3.6, pp.~321 and 327--328]{DvornicichZannier2001}.
The relative profinite formalism is developed by
\citet[\S2.1, Proposition~2.4]{WilkesRelative2019}, and its extension and
projectivity interpretation by \citet{WilkesRelative2025}.
The following explicit example verifies the additional generation and
invariant hypotheses which arise in the present source argument.

\begin{proposition}\label{prop:conjugate-source-not-enough}
There exist a profinite group \(\Gamma=P\rtimes\Delta\), with \(P\) a
finitely generated free pro-three group and \(|\Delta|\) prime to three,
a closed subgroup \(D\), and a finite free continuous
\(\mathbf Z_3[\Gamma]\)-module \(T\), such that:
\begin{enumerate}
\item the image of \(P\) on \(T\) is pro-three;
\item the conjugates of \(D\) generate \(\Gamma\);
\item \(T^D=0\);
\item \(H^1(\Gamma,T)\longrightarrow H^1(D,T)\) is not injective,
even after tensoring with \(\mathbf Q_3\).
\end{enumerate}
The same nonzero class restricts to zero on every conjugate of \(D\).
\end{proposition}
\begin{proof}
Index five symbols by \(i\in\mathbf Z/5\mathbf Z\). Let \(P\) be free
pro-three on \(x_i\), and let the dihedral group \(\Delta=D_{10}\) act
by its usual permutations of the symbols. Put
\[
 T=\left\{(a_i)\in\mathbf Z_3^5:\sum_i a_i=0\right\}.
\]
Let \(\Delta\) permute the coordinates. The homomorphism
\(\epsilon:P\to\mathbf Z_3\), defined by \(\epsilon(x_i)=1\), is
\(\Delta\)-invariant. Let \(g\in P\) act by \(4^{\epsilon(g)}\).
This gives a continuous action of \(\Gamma=P\rtimes\Delta\) on \(T\);
its \(P\)-image is pro-three.

Let \(\Delta_i\) be the reflection subgroup fixing \(i\), and put
\[
 D_i=\langle x_i\rangle\rtimes\Delta_i,\qquad
 b_i=5e_i-\sum_j e_j .
\]
Then \(b_i\in T^{\Delta_i}\), and \(\sigma b_i=b_{\sigma i}\).
Assigning \(x_i\) to \((3b_i,4)\) in the pro-three semidirect product
\(T\rtimes\overline{\langle4\rangle}\) gives, by freeness, a continuous
cocycle \(c:P\to T\) with \(c(x_i)=3b_i\). It is
\(\Delta\)-equivariant, so it extends to \(\Gamma\) with \(c|_\Delta=0\).

On \(D_i\), the cocycle is the coboundary of \(b_i\): it takes the value
\((4-1)b_i=3b_i\) on \(x_i\) and zero on \(\Delta_i\).
If it were a global coboundary with parameter \(b\in T\), its zero
restriction to \(\Delta\) would imply \(b\in T^\Delta=0\), a
contradiction. The same argument applies over \(\mathbf Q_3\).

Multiplication by \(4-1=3\) is injective on \(T\), so \(T^{D_i}=0\).
The subgroups \(D_i\) generate \(P\) through the \(x_i\), and generate
\(\Delta\) through the five reflections. Hence they generate \(\Gamma\).
They are conjugate under \(\Delta\). Take \(D=D_0\).
Restriction to every further \(\Gamma\)-conjugate is also zero, by
the conjugacy compatibility of restriction in cohomology.
\end{proof}

\begin{remark}\label{rem:conjugate-source-scope}
This is an abstract continuous group-cohomology counterexample, not an
arithmetic Selmer-section counterexample. It shows that equivariance,
generation by conjugate local subgroups, and zero local invariants do
not justify replacing a single surjective local source by its normal
closure. The local coboundary parameters still require compatibility.
\end{remark}

\begin{lemma}[A module-specific sufficient condition]
\label{lem:source-local-parameter}
Let \(\Gamma=P\rtimes\Delta\), with \(P\) pro-three and
\(\Delta\) finite of order prime to three. Let
\(D=P_0\rtimes\Delta_0\) be a closed subgroup, where
\(\Delta_0\subset\Delta\), and suppose the \(\Delta\)-conjugates of
\(P_0\) generate \(P\). If \(T\) is finite free over \(\mathbf Z_3\)
and \(T^{P_0}=0\), there is a natural injection
\[
 \ker\!\left(H^1(\Gamma,T)\to H^1(D,T)\right)
       \hookrightarrow T^{\Delta_0}/T^\Delta.
\]
In particular equality of the two invariant lattices implies
localization injectivity.
\end{lemma}
\begin{proof}
Normalize the cocycle to vanish on \(\Delta\), using
\(H^1(\Delta,T)=0\). A local coboundary then has a unique parameter
\(b\), since \(T^{P_0}=0\), and \(b\in T^{\Delta_0}\).
Changing the normalized cocycle changes \(b\) by an element of
\(T^\Delta\). If its class in the displayed quotient is zero,
subtract that global coboundary. The cocycle now vanishes on
\(\Delta\) and \(P_0\), hence on every \(\Delta\)-conjugate of \(P_0\)
and therefore on \(\Gamma\).
\end{proof}
\subsection{The relation and character conditions for joint localization}

The preceding counterexample leaves a precise compatibility condition.
It can be expressed either through relations among local source groups
or, for abelian coefficient actions, through selected character components
of an arithmetic Galois module.

\begin{proposition}[The exact local-parameter condition]
\label{prop:source-exact-relation-condition}
Let \(\Gamma=P\rtimes\Delta\), where \(P\) is pro-three and
\(\Delta\) is finite of order prime to three. Let
\(D=P_0\rtimes\Delta_0\), with \(\Delta_0\subset\Delta\), be a
closed subgroup. Suppose the \(\Delta\)-conjugates of \(P_0\)
generate \(P\). Let \(T\) be a finite free continuous
\(\mathbf Z_3[\Gamma]\)-module with \(T^{P_0}=0\).

Form the free pro-three product \(F\) of copies of \(P_0\) indexed
by \(\Delta/\Delta_0\), with the action induced by the conjugation
action of \(\Delta_0\) on \(P_0\). The conjugation maps give a
surjection \(F\to P\). Put \(R=\ker(F\to P)\).
For \(b\in T^{\Delta_0}\), let \(c_b\) be the cocycle on \(F\)
whose restriction to the factor indexed by \(\sigma\Delta_0\)
is the coboundary with parameter \(\sigma b\).
Then
\begin{equation}
 \ker\!\left(H^1(\Gamma,T)\longrightarrow H^1(D,T)\right)
 \simeq
 \frac{\{b\in T^{\Delta_0}:c_b|_R=0\}}{T^\Delta}.
 \label{eq:source-exact-relation-condition}
\end{equation}
\end{proposition}
\begin{proof}
The definition of \(c_b\) is independent of the chosen coset
representative because \(b\) is fixed by \(\Delta_0\).
The action of \(F\) on \(T\) factors through \(P\), whose image
is pro-three. The universal property of the free pro-three product,
applied to the cocycle graph in \(T\rtimes\rho(P)\), gives a
continuous cocycle \(c_b\). It is \(\Delta\)-equivariant.
Since \(R\) acts trivially on \(T\), its restriction to \(R\)
is a homomorphism, and it descends to \(P\) precisely when this
homomorphism is zero. In that case it extends to \(\Gamma\)
by taking value zero on \(\Delta\), and it is a coboundary
on \(D\).

Conversely, a class in the displayed kernel has a representative
vanishing on \(\Delta\), because \(|\Delta|\) is a unit in
\(\mathbf Z_3\). Its unique local coboundary parameter lies in
\(T^{\Delta_0}\); uniqueness uses \(T^{P_0}=0\).
Equivariance determines its restrictions to every factor of \(F\).
The existence of the original cocycle on \(P\) is exactly
\(c_b|_R=0\). Two normalized representatives differ by a
coboundary with parameter in \(T^\Delta\). This proves
\eqref{eq:source-exact-relation-condition}.
\end{proof}

\paragraph{An arithmetic module testing only the required characters.}
The compact kernel of localization at the places over the coefficient
prime is a familiar fine-Selmer object; compare
\citet[Eq.~(1.2), author p.~2]{WuthrichFineSelmer2007}.
The following formulation gives a direct character condition and
does not assert vanishing of that kernel.

\begin{proposition}[An abelian-character localization condition]
\label{prop:cm-character-localization-condition}
Let \(E/\mathbf Q\) be finite Galois. Let \(S\) be a finite
set of rational primes containing three and the ramification
of \(E/\mathbf Q\) and of a finite-dimensional continuous
\(\mathbf Q_3\)-representation \(V\).
Let \(E_S\) be the maximal extension of \(E\) unramified outside
the primes over \(S\). Suppose \(E_\infty\subset E_S\) is Galois
over \(\mathbf Q\), is abelian over \(E\), and trivializes
the action of \(G_E\) on \(V\). Put
\[
 H=\operatorname{Gal}(E_\infty/E),\qquad
 \Delta=\operatorname{Gal}(E/\mathbf Q).
\]
Suppose that, over a finite coefficient extension
\(\Omega/\mathbf Q_3\), the \(H\)-representation
\(V_\Omega=V\otimes_{\mathbf Q_3}\Omega\) is a direct sum
of characters, each nontrivial on \(H\) and on every local
decomposition image \(H_v\) for \(v\mid3\).

Let \(X_S^{[3]}(E_\infty)\) be the maximal abelian pro-three
quotient of \(\operatorname{Gal}(E_S/E_\infty)\) in which
the decomposition groups at every prime above three vanish.
Thus it classifies abelian pro-three extensions unramified
outside \(S\) and split at every three-adic prime.
There is a natural isomorphism
\begin{equation}
\begin{split}
 &\ker\!\left(H^1(G_{\mathbf Q,S},V)
            \longrightarrow H^1(\mathbf Q_3,V)\right)
                \otimes_{\mathbf Q_3}\Omega\\
 &\qquad\simeq
 \operatorname{Hom}^{\mathrm{cts}}_H
       \bigl(X_S^{[3]}(E_\infty),V_\Omega\bigr)^\Delta.
\end{split}
\label{eq:cm-character-localization-condition}
\end{equation}
In particular only the character components which occur in
\(V_\Omega\) need to vanish.
\end{proposition}
\begin{proof}
Let \(\chi\) be one of the characters. Choose \(h\in H\)
with \(\chi(h)\ne1\). The inner action of \(h\) on continuous
group cohomology is trivial. Since \(H\) is abelian, the same
action is multiplication by \(\chi(h)\) on
\(H^j(H,\Omega(\chi))\). The scalar \(\chi(h)-1\) is invertible,
so these cohomology groups vanish for all \(j\).
Inflation--restriction therefore identifies
\[
 H^1(G_{E,S},\Omega(\chi))
 \simeq
 \operatorname{Hom}^{\mathrm{cts}}_H
 \bigl(\operatorname{Gal}(E_S/E_\infty)^{\mathrm{ab},(3)},
       \Omega(\chi)\bigr).
\]
Continuous additive images in the coefficient line are pro-three,
which explains the indicated abelian pro-three quotient.

The same argument applies locally because
\(\chi|_{H_v}\ne1\). Consequently vanishing of a local class
is equivalent to the corresponding homomorphism killing the
decomposition subgroup over that place. Summing over the
characters gives the joint-localization kernel over \(E\).

Finite-group cohomology with coefficients in \(\Omega\) vanishes
in positive degrees by averaging. Globally, inflation--restriction
through \(E/\mathbf Q\) identifies the rational cohomology
with the \(\Delta\)-invariants over \(E\). Locally, its analogous
form identifies \(H^1(\mathbf Q_3,V_\Omega)\) with the
\(\Delta\)-invariants of the product over \(v\mid3\).
Taking kernels proves \eqref{eq:cm-character-localization-condition}.
The action of \(\Delta\) on the right is well defined: changing
a lift in \(\operatorname{Gal}(E_\infty/\mathbf Q)\) changes
both source and target by the same \(H\)-action.
\end{proof}

\begin{lemma}[Integral recovery from rational injectivity]
\label{lem:localization-integral-recovery}
Let \(T\) be a stable finite free \(\mathbf Z_3\)-lattice in \(V\),
and assume \(V^{G_{\mathbf Q_3}}=0\).
If
\[
 H^1(G_{\mathbf Q,S},V)\longrightarrow H^1(\mathbf Q_3,V)
\]
is injective, then so is the corresponding map with coefficients
in \(T\).
\end{lemma}
\begin{proof}
A class in the integral kernel is zero after passing to \(V\).
The exact sequence \(0\to T\to V\to V/T\to0\) makes it the
boundary of an element of \((V/T)^{G_{\mathbf Q,S}}\).
The local boundary map from \((V/T)^{G_{\mathbf Q_3}}\) is
injective because \(V^{G_{\mathbf Q_3}}=0\).
The global invariant is also a local invariant, and its local
boundary is zero; it is therefore zero. This proves the claim.
\end{proof}

\begin{lemma}[Local nontriviality for the eleven coefficient characters]
\label{lem:eleven-coefficient-local-nontriviality}
Let \(N_{11}\) be the geometric pro-three group of \(U_{11}\),
and put \(E=\mathbf Q(\zeta_{11})\). Over a finite splitting
coefficient field, every character in a positive ordinary
lower-central grade of \(N_{11}\) is nontrivial on
\(G_{E_v}\) for each \(v\mid3\). The same holds for every
nonzero character in a stable quotient of such a grade.
In particular these rational coefficient modules have zero
local invariants at three.
\end{lemma}
\begin{proof}
The decomposition \(J_{11}\sim A_{11}^3\times B^6\), with
\(B=\operatorname{Jac}(y^{11}=x(1-x)^2)\), reduces the
compact degree-one calculation to two CM types. The
holomorphic types are
\[
 H_A=\{1,2,3,4,5\},\qquad H_B=\{1,2,3,6,7\};
\]
equivalently the reflex types are
\[
 R_A=\{1,6,4,3,9\},\qquad R_B=\{1,6,4,2,8\}.
\]
These follow from the Fermat quotient type formula
\citep[\S2.1, author p.~5]{ShuExpanded2021}.
The two Frobenius-three orbits are
\[
 C=\langle3\rangle=\{1,3,9,5,4\},\qquad -C.
\]
The intersections of \(R_A\) with these orbits have sizes
four and one; those of \(R_B\) have sizes two and three.
Multiplication by a coefficient automorphism either preserves
or exchanges the two orbits, so every conjugate type meets
both.

At the good prime three, local CM reciprocity on an open
group of principal units is the inverse reflex norm
\citep[Theorem~5.1, pp.~483--484]{Lombardo2017}.
Every compact character therefore has, at every \(v\mid3\),
a nonempty set of algebraic unit exponents, all of the same
sign. The boundary characters are cyclotomic and have the
same sign. A character in a positive tensor power is a
product with nonnegative multiplicities, so its unit exponent
vector cannot be zero. Indeed, after taking the three-adic
logarithm on a sufficiently small unit group, the distinct
embeddings are linearly independent and a nonzero exponent
vector gives a nonzero linear map.

Ordinary lower-central grades are generated by brackets in
degree one, and their actions are subquotients of these tensor
powers. Coordinate-cusp differences are killed by eleven, so the
generalized-Jacobian extension splits on three-adic Tate modules,
as explained in Remark~\ref{rem:boundary-finite-invariant-obstruction}.
Over the splitting coefficient field the action is therefore
diagonal, so subquotients retain only characters from that
list. This proves the claim.
\end{proof}

\begin{remark}[A degree-two boundary component]
\label{rem:boundary-finite-invariant-obstruction}
For \(a=11\) or \(13\), put \(E=\mathbf Q(\zeta_a)\).
The boundary module of \(U_a\) has the integral decomposition
\[
 B_a\simeq\mathbf Z_3(1)^2\oplus
     \left(\operatorname{Ind}_{G_E}^{G_{\mathbf Q}}
                       \mathbf Z_3(1)\right)^3.
\]
There are three rational boundary points and three primitive
cyclotomic orbits; eliminating one rational basis vector proves
the decomposition.
Coordinate-cusp differences are killed by \(a\), so the
generalized-Jacobian extension splits on three-adic Tate modules.
Hence \(B_a\) is a summand of the geometric abelian module.
For the free geometric pro-three group, its second lower-central
grade is the exterior square of that module.

Choose a nontrivial character \(\eta\) of
\((\mathbf Z/a)^\times/\langle3\rangle\), using a coefficient
extension when necessary. Bracketing a rational boundary line
with the \(\eta\)-line in an induced summand gives
\(\Omega(2)\otimes\eta\) in the second grade.
For eleven one can take the odd quadratic character; for thirteen
the displayed quotient is \(C_4\).
The finite character \(\eta\) is trivial on the local group at
three and nontrivial globally. The square of the finite
mod-three cyclotomic character is trivial.
Thus the associated finite coefficient representation has more
local invariants than global invariants already in degree two.
An invariant-equality criterion alone cannot handle all these
grades. This does not establish a nonzero localization kernel:
the relation condition or the arithmetic character condition
must still be tested.
\end{remark}

\begin{remark}[Finite ramification and the base-exponent problem]
\label{rem:joint-localization-ramification-scope}
The character condition concerns classes in
\(H^1(G_{\mathbf Q,S},V)\). It does not force an arbitrary
affine Selmer section to have its whole relative cocycle
unramified outside \(S\).
At a prime exponent different from three, compact
three-adic Kummer classes do not detect differences between
coordinate cusps, since those differences have order prime to
three. Compact localization and one locally cuspidal
representative therefore do not by themselves restrict boundary
ramification to one global cusp direction.

If the global abelian class is already that of one cusp,
the point-or-cusp local condition can control inertia in the
subsequent central induction. A nonzero boundary inertia vector
forces the same cusp direction; a zero vector makes the whole
inertia trivial. This is a separate step, preceding application
of a finite-ramification cohomological condition.
The propositions above consequently supply formal
module-specific conditions, not a new base-eleven or
base-thirteen pro-three cuspidality theorem.
\end{remark}
